% !TeX root = ../thesis.tex
\chapter{Giới thiệu đề tài}
\label{chap:intro}

\section{Bối cảnh và bài toán tìm kiếm câu hỏi}
Trong các cộng đồng hỏi đáp, người dùng thường gặp một vấn đề đã được thảo luận trước đó nhưng không biết câu hỏi cũ sử dụng cách diễn đạt nào. Một người có thể mô tả hiện tượng bằng ngôn ngữ tự nhiên, trong khi câu hỏi liên quan lại đề cập tên công cụ, thông báo lỗi hoặc khái niệm chuyên môn. Hệ thống tìm kiếm cần đưa những câu hỏi phù hợp lên đầu danh sách để người dùng tận dụng tri thức đã có. Khó khăn không chỉ nằm ở việc tìm các văn bản có chung từ ngữ mà còn ở việc nhận biết những mô tả khác nhau của cùng một nhu cầu.

Trong bối cảnh kỹ thuật, sự khác biệt về diễn đạt thường tồn tại đồng thời với yêu cầu khớp chính xác. Chẳng hạn, một truy vấn hỏi cách ``thực thi chương trình theo lịch'' có thể liên quan đến tài liệu nói về công cụ lập lịch, dù hai văn bản không dùng cùng cụm từ. Ngược lại, nếu truy vấn chứa một mã lỗi hoặc tùy chọn dòng lệnh cụ thể, việc bỏ qua chi tiết đó có thể làm thay đổi hoàn toàn vấn đề cần giải quyết. Hai tình huống này là ví dụ minh họa cho nhu cầu kết hợp tín hiệu từ khóa và tín hiệu ngữ nghĩa; chúng không phải các trường hợp thực nghiệm được dùng để tính kết quả trong báo cáo.

Đề tài xem xét bài toán truy xuất câu hỏi. Đầu vào của hệ thống là một truy vấn dạng văn bản và một kho câu hỏi đã được lập chỉ mục. Đầu ra là danh sách câu hỏi được sắp xếp theo mức độ liên quan. Nhãn trong bộ dữ liệu cho biết những câu hỏi nào có quan hệ phù hợp với truy vấn để đánh giá danh sách đó. Hệ thống không sinh một câu trả lời mới và không đánh giá khả năng diễn đạt của mô hình ngôn ngữ. Việc xác định phạm vi như vậy giúp tập trung vào quyết định truy xuất và chi phí cần thiết để tạo ra danh sách kết quả.

CQADupStack cung cấp môi trường đánh giá phù hợp với bài toán này nhờ các câu hỏi và quan hệ trùng lặp từ nhiều cộng đồng Stack Exchange \cite{cqadupstack2015}. Quá trình nghiên cứu bắt đầu với Programmers và Unix, sau đó mở rộng đến các chủ đề khác. Vì tập xác nhận cuối còn bao gồm English và Gaming, báo cáo sử dụng tên gọi chung là tìm kiếm câu hỏi trên CQADupStack. Những nhận xét rút ra từ một miền kỹ thuật không mặc nhiên áp dụng cho mọi miền của tập dữ liệu.

\section{Vấn đề chất lượng và chi phí truy xuất}
Truy xuất theo từ khóa thường có chi phí thấp và khai thác tốt sự xuất hiện của các thuật ngữ cụ thể. Tuy nhiên, kết quả có thể hạn chế khi truy vấn và câu hỏi liên quan biểu đạt cùng ý bằng những từ khác nhau. Truy xuất Dense sử dụng biểu diễn ngữ nghĩa để bổ sung một nguồn bằng chứng khác về mức độ liên quan. Khi kết hợp hai nguồn, hệ thống Hybrid có thể cải thiện danh sách kết quả so với chỉ dùng một phương pháp. Các kết quả benchmark như BEIR đồng thời cho thấy cần kiểm chứng khả năng tổng quát hóa và giữ những đối chứng từ khóa mạnh khi đánh giá phương pháp truy xuất \cite{beir2021}.

Lợi ích chất lượng của Hybrid đi kèm phần công việc phát sinh cho mỗi truy vấn mới. Vector tài liệu có thể được tính và lưu trước, nhưng truy vấn cần được mã hóa tại thời điểm phục vụ. Sau bước này, hệ thống còn phải tìm kiếm vector và kết hợp kết quả với danh sách BM25. Vì vậy, kể cả khi chỉ mục đã sẵn sàng, luôn chạy Hybrid vẫn phải trả chi phí cho nhánh Dense ở mọi yêu cầu. Báo cáo quan tâm đến chi phí xử lý từng truy vấn trong điều kiện chỉ mục và mô hình đã được tải, thay vì đồng nhất chi phí này với thời gian xây dựng hệ thống ban đầu.

Vấn đề nghiên cứu xuất hiện khi mức độ cần thiết của Dense thay đổi theo truy vấn. Nếu BM25 đã đưa câu hỏi liên quan lên đầu danh sách, việc tiếp tục chạy Dense có thể chỉ đem lại lợi ích nhỏ hoặc không làm thay đổi chất lượng. Ngược lại, nếu BM25 thất bại, bỏ qua Dense có thể làm mất cơ hội tìm được kết quả phù hợp. Một chính sách áp dụng giống nhau cho mọi truy vấn không tận dụng sự khác biệt đó. Nhu cầu đặt ra là phân bổ chi phí Dense cho những truy vấn có khả năng cần đến nó, thay vì chỉ giảm tỷ lệ sử dụng Dense một cách tùy ý.

Việc lựa chọn cũng có chi phí riêng. Bộ quyết định càng phức tạp thì phần tiết kiệm được từ bỏ qua Dense càng có nguy cơ bị triệt tiêu. Một bộ quyết định phải mã hóa truy vấn trước khi lựa chọn còn có thể đã trả phần lớn chi phí mà nó muốn tránh. Do đó, đề tài đặt quyết định sau BM25 nhưng trước encoder. Tại thời điểm này, hệ thống đã có một danh sách từ khóa và những thống kê liên quan, nhưng chưa có biểu diễn ngữ nghĩa của truy vấn mới.

Giảm số lần chạy Dense chưa đủ để khẳng định hệ thống tốt hơn trong sử dụng thực tế. Quyết định sai có thể làm giảm chất lượng, còn phần xử lý bổ sung có thể làm mức giảm thời gian nhỏ hơn tỷ lệ bỏ qua Dense. Độ trễ của các yêu cầu chậm nhất cũng có thể không cải thiện nếu những truy vấn đó vẫn đi qua nhánh đắt. Bài toán vì thế cần được đánh giá đồng thời theo chất lượng, tỷ lệ mã hóa và chi phí đầu cuối, với điều kiện so sánh được kiểm soát.

\section{Giả thiết và câu hỏi nghiên cứu}
Các nghiên cứu về độ khó truy vấn cho thấy hiệu quả truy xuất phụ thuộc quan hệ giữa nhu cầu thông tin và dữ liệu \cite{querydifficulty2006}. Giả thiết của đề tài là thông tin từ truy vấn và kết quả BM25 có thể cung cấp tín hiệu về khả năng BM25 đáp ứng nhu cầu tìm kiếm. Ví dụ, khoảng cách điểm giữa các kết quả đầu hoặc mức độ phù hợp của từ ngữ trong truy vấn với kho dữ liệu có thể phản ánh sự chắc chắn của kết quả từ khóa. Nếu những tín hiệu này đủ hữu ích, một mô hình nhẹ có thể giúp ưu tiên sử dụng Dense mà không cần nhìn vào kết quả Dense trước đó. Đây là giả thiết cần kiểm chứng bằng thực nghiệm, không phải bảo đảm rằng mọi truy vấn khó đối với BM25 đều được Dense giải quyết.

Từ giả thiết đó, báo cáo đặt ra bốn câu hỏi nghiên cứu có quan hệ với nhau. Câu hỏi thứ nhất là liệu bộ quyết định học được có mang lại chất lượng tốt hơn một quy tắc đơn giản và lựa chọn ngẫu nhiên khi các phương pháp được sử dụng cùng số lượt Dense. Câu hỏi này tách giá trị của việc lựa chọn truy vấn khỏi lợi ích chỉ đến từ việc dùng thêm tài nguyên.

Câu hỏi thứ hai là liệu PE-GH có thể giảm số lượt mã hóa trong khi chất lượng vẫn nằm trong biên suy giảm cho phép so với luôn chạy Hybrid. Thay vì yêu cầu hai phương pháp đạt điểm hoàn toàn giống nhau, đề tài xác định trước một mức mất mát tối đa để đánh giá ý nghĩa của việc tiết kiệm chi phí. Biên đánh giá này phải được chốt trước khi xem kết quả xác nhận để tránh điều chỉnh tiêu chí theo số liệu thu được.

Câu hỏi thứ ba là việc giảm số lượt Dense có chuyển thành giảm độ trễ và CPU khi hệ thống phục vụ truy vấn thật hay không. Câu trả lời cần dựa trên phép đo trực tiếp, bởi độ trễ còn phụ thuộc bộ quyết định, nhánh chi phí thấp và chi phí kết hợp kết quả. Báo cáo đồng thời xem xét p95 và bộ nhớ để xác định những khía cạnh chưa được cải thiện.

Câu hỏi thứ tư là hệ thống cuối cải thiện đến đâu so với phiên bản ban đầu và những giới hạn nào còn tồn tại. So sánh này bao gồm thay đổi encoder, nhánh truy xuất và cách học của bộ quyết định, nên phản ánh tác động của cả hệ thống. Để tránh quy kết quá mức, báo cáo phân biệt câu hỏi về nâng cấp toàn hệ thống với câu hỏi về giá trị riêng của bộ quyết định.

\section{Mục tiêu, phạm vi và đóng góp}
Mục tiêu phương pháp của đề tài là xây dựng một chính sách truy xuất lai có chọn lọc, trong đó truy vấn mới chỉ được mã hóa khi bộ quyết định lựa chọn nhánh Dense. Hệ thống cuối PE-GH v2 dự đoán khả năng BM25 không tìm được kết quả liên quan, dựa trên thông tin có sẵn trước khi chạy Dense. Ngân sách truy xuất cho phép điều chỉnh mức sử dụng tài nguyên để quan sát những điểm đánh đổi khác nhau giữa chất lượng và chi phí.

Mục tiêu thực nghiệm là đánh giá chính sách này bằng các phép so sánh đủ mạnh. Ngoài BM25 và luôn chạy Hybrid, đề tài giữ quy tắc BM25-margin, lựa chọn ngẫu nhiên và đối chứng biết trước chất lượng của từng nhánh để xác định mức lợi ích thực sự của việc học. Mô hình được phát triển trên sáu miền, sau đó giữ cố định để xác nhận trên sáu miền khác. Những giả thuyết chính, tiêu chí lựa chọn và biên chất lượng được đăng ký trước giai đoạn xác nhận trong quy trình thực nghiệm của nhóm.

Đóng góp thứ nhất là cách tổ chức một bộ quyết định nhẹ theo thông tin mà hệ thống thực sự quan sát được trước encoder, cùng việc đánh giá giới hạn của lựa chọn đó. Ý tưởng lựa chọn chiến lược truy xuất đã có trong nghiên cứu trước \cite{arabzadeh2021}; điểm cần chứng minh trong đề tài là hiệu quả của cấu hình cụ thể trên bài toán, dữ liệu và điều kiện chi phí đã xác định. Báo cáo không xem việc dùng một bộ quyết định nhẹ hoặc đặt tên phương pháp là bằng chứng về tính mới.

Đóng góp thứ hai là bộ bằng chứng thực nghiệm nối chất lượng với chi phí đo thật. So sánh ở cùng ngân sách cho biết việc lựa chọn có hữu ích hay không, còn phép đo đầu cuối cho biết phần lợi ích đó có ý nghĩa vận hành đến đâu. Khoảng tin cậy, kết quả từng miền và các trường hợp không thuận lợi được giữ trong diễn giải. Đóng góp này giúp xác định điều kiện sử dụng hợp lý của phương pháp thay vì chỉ công bố một tỷ lệ tăng tốc tổng hợp.

Phạm vi kết luận được giới hạn ở CQADupStack, mô hình truy xuất ngữ nghĩa đã chọn và môi trường CPU đã đo. Đánh giá chất lượng xác nhận sử dụng sáu miền; đo khi phục vụ truy vấn sử dụng Programmers và Unix với bộ quyết định được huấn luyện không có miền đang đo. Vì vậy, báo cáo không suy rộng số đo thời gian sang mọi chủ đề, phần cứng hoặc tải đồng thời. Kết luận chỉ áp dụng cho phiên bản đã được đánh giá, không tự động áp dụng cho những thay đổi được thực hiện sau đó.

\section{Bố cục báo cáo}
Chương~\ref{chap:background} giới thiệu ngắn những khái niệm cần thiết và phân tích các hướng nghiên cứu gần. Chương~\ref{chap:method} trình bày nguyên lý PE-GH và cơ sở lựa chọn phương pháp. Chương~\ref{chap:experiment} mô tả dữ liệu, đối chứng và quy trình kiểm chứng. Chương~\ref{chap:results} phân tích bằng chứng về chất lượng, chi phí và giới hạn. Cuối cùng, Chương~\ref{chap:conclusion} trả lời các câu hỏi nghiên cứu và xác định hướng phát triển phù hợp với những gì đã được chứng minh.
